\documentclass[11pt]{article}
\usepackage[a4paper,margin=0.9in]{geometry}
\usepackage[T1]{fontenc}
\usepackage{helvet}
\renewcommand{\familydefault}{\sfdefault}
\usepackage{xcolor}
\usepackage{booktabs}
\usepackage{array}
\usepackage{colortbl}
\usepackage{float}
\usepackage{amsmath}
\usepackage[font=footnotesize,labelformat=empty]{caption}
\usepackage{titlesec}
\usepackage{fancyhdr}
\usepackage{enumitem}
\usepackage{lastpage}
\usepackage[hidelinks]{hyperref}

\definecolor{bcnavy}{RGB}{11,31,58}
\definecolor{bcblue}{RGB}{32,74,135}
\definecolor{bcgold}{RGB}{176,141,87}
\definecolor{bcgrey}{RGB}{90,98,112}
\definecolor{bcrule}{RGB}{210,214,220}
\definecolor{bckeep}{RGB}{232,237,244}
\definecolor{bcact}{RGB}{247,241,228}
\definecolor{bcwarn}{RGB}{250,238,238}

\titleformat{\section}{\color{bcnavy}\Large\bfseries}{\thesection}{0.6em}{}
\titleformat{\subsection}{\color{bcblue}\large\bfseries}{\thesubsection}{0.6em}{}
\titlespacing*{\section}{0pt}{1.3em}{0.45em}

\pagestyle{fancy}\fancyhf{}
\renewcommand{\headrulewidth}{0pt}
\renewcommand{\footrulewidth}{0.4pt}
\renewcommand{\footrule}{\hbox to\headwidth{\color{bcrule}\leaders\hrule height \footrulewidth\hfill}}
\fancyfoot[L]{\footnotesize\color{bcgrey}Bayesian Capital}
\fancyfoot[C]{\footnotesize\color{bcgrey}Confidential --- Internal Research}
\fancyfoot[R]{\footnotesize\color{bcgrey}Page \thepage\ of \pageref{LastPage}}

\setlength{\parindent}{0pt}\setlength{\parskip}{0.5em}

\begin{document}

\begin{titlepage}
\vspace*{1.2cm}
{\color{bcgold}\rule{\textwidth}{2pt}}\\[0.4cm]
{\color{bcnavy}\Huge\bfseries Seven-Name Book}\\[0.25cm]
{\color{bcnavy}\Huge\bfseries Implementation Specification}\\[0.35cm]
{\color{bcblue}\LARGE Daily Hybrid and Weekly Monday Tranche}\\[0.5cm]
{\color{bcgold}\rule{\textwidth}{1pt}}\\[0.6cm]
{\large\color{bcgrey} Five names on the daily Bayesian / Ornstein--Uhlenbeck model}\\[0.2cm]
{\large\color{bcgrey} Two names on the weekly single-tranche model}\\[0.2cm]
{\large\color{bcgrey} Every name above the 50\% floor in the tested half of the sample}\\[1.2cm]
{\color{bcnavy}\large\bfseries Bayesian Capital}\\[0.15cm]
{\color{bcgrey} Quantitative Research}\\[0.15cm]
{\color{bcgrey} 4 August 2026}\\[1.6cm]

\fbox{\parbox{0.92\textwidth}{\color{bcnavy}\textbf{Purpose.}\\ \color{black}
The operative configuration for a book in which every constituent clears a 50\% annualised
floor. Four names --- TSLA, SPOT, SOFI and PLTR --- are removed rather than repaired, having
failed on both models. NVDA and AVGO move from the daily model to a weekly single Monday
tranche. The remaining five keep their existing daily parameters unchanged.

All figures are on the execution-verified basis: every same-day round trip is checked against
minute or 5-minute bars, so no fill depends on an intraday ordering that cannot be known from
daily data.}}
\end{titlepage}

\section{The book}

\begin{table}[H]\centering\small
\begin{tabular}{llrrrr}
\toprule
\textbf{Name} & \textbf{Model} & \textbf{Verified} & \textbf{Tested half} & \textbf{Planning} &
\textbf{Weight}\\
\midrule
\rowcolor{bckeep} RKLB & Daily hybrid   & 111\% & 65\%  & 115\% & 14.29\%\\
TSM  & Daily hybrid   & 68\%  & 89\%  & 55\%  & 14.29\%\\
\rowcolor{bckeep} VST  & Daily hybrid   & 65\%  & 55\%  & 57\%  & 14.29\%\\
VRT  & Daily hybrid   & 65\%  & 168\% & 53\%  & 14.29\%\\
\rowcolor{bckeep} MU   & Daily hybrid   & 54\%  & 160\% & 54\%  & 14.29\%\\
AVGO & Weekly Monday  & 91\%  & 96\%  & 60\%  & 14.29\%\\
\rowcolor{bckeep} NVDA & Weekly Monday  & 82\%  & 82\%  & 58\%  & 14.29\%\\
\midrule
\textbf{Book} & equal weight & \textbf{77\%} & --- & \textbf{65\%} & 100\%\\
\textbf{Book ex-RKLB} & & \textbf{71\%} & --- & \textbf{56\%} & \\
\bottomrule
\end{tabular}
\caption{Verified = full sample at the deployed configuration. Tested half = weeks 60--120,
the portion scored outside every training window. Planning = the figure to budget on, derived
in \S5. Daily names annualised on the workbook's 2.2-year convention; weekly names on the true
2.31-year span.}
\end{table}

The planning column is materially below the verified column for every name except RKLB and MU.
That gap is deliberate and is the subject of \S5 --- it is the difference between what the
configuration scored and what survived validation.

\section{Daily names --- parameters}

Five names, unchanged from the current workbooks. Nine attempts to improve these parameters were
tested out of sample across this research and all were rejected; they are reproduced here exactly
as they stand. Cell references are to the \texttt{2 Tranche StopLoss 50d} sheet.

\begin{table}[H]\centering\footnotesize
\setlength{\tabcolsep}{4pt}
\begin{tabular}{lrrrrrr}
\toprule
& \multicolumn{6}{c}{\textbf{Bayes sleeve}}\\
\cmidrule(l){2-7}
\textbf{Name} & $\lambda$ (B2) & $\phi_L$ (D2) & $\psi$ (F2) & $k$ (H2) & prem (J2) &
peak cap (L2)\\
\midrule
\rowcolor{bckeep} RKLB & 0.352847 & 0.237406 & 0.113048 & 0.712293 & 0.026842 & 0.008015\\
TSM  & 0.462001 & 0.254287 & 0.065039 & 1.004530 & 0.015652 & 0.030484\\
\rowcolor{bckeep} VST  & 0.373520 & 0.264322 & 0.096187 & 1.385660 & 0.021816 & 0.013663\\
VRT  & 0.420699 & 0.425851 & 0.081656 & 1.116400 & 0.021691 & 0.048379\\
\rowcolor{bckeep} MU   & 0.600000 & 0.263600 & 0.010000 & 1.054800 & 0.024300 & 0.038000\\
\bottomrule
\end{tabular}
\end{table}

\begin{table}[H]\centering\footnotesize
\setlength{\tabcolsep}{5pt}
\begin{tabular}{lrrrrrr}
\toprule
& \multicolumn{4}{c}{\textbf{OU sleeve}} & \multicolumn{2}{c}{\textbf{Common}}\\
\cmidrule(lr){2-5}\cmidrule(l){6-7}
\textbf{Name} & $W$ (B3) & buffer $k$ (D3) & prem (H3) & cap (J3) & stop (T2) & comm (N2)\\
\midrule
\rowcolor{bckeep} RKLB & 85  & 0.220593 & 0.022147 & 0.033647 & 50 & 0.005\\
TSM  & 77  & 0.399229 & 0.010022 & 0.020893 & 50 & 0.005\\
\rowcolor{bckeep} VST  & 122 & 0.238599 & 0.014498 & 0.052746 & 50 & 0.005\\
VRT  & 80  & 0.567888 & 0.026603 & 0.059499 & 50 & 0.005\\
\rowcolor{bckeep} MU   & 91  & 0.300000 & 0.025000 & 0.029200 & 50 & 0.005\\
\bottomrule
\end{tabular}
\caption{Interest on idle cash (R2) is 3.14\% for all names. The stop period is a uniform 50
calendar days: a variable stop was tested, gained 8.9pp in sample and failed walk-forward at
13 folds of 30.}
\end{table}

\subsection{Sleeve split}

Bayes share \textbf{0.75} for all five names, entered flat rather than per name. Per-name Bayes
shares were tested and not supported: training windows selected a Bayes-heavy share for
essentially every name, so the per-name freedom bought nothing that the flat tilt did not
already provide.

\section{Weekly names --- parameters}

NVDA and AVGO run a \textbf{single Monday tranche} each. The three-tranche staggered version was
tested and the simplification costs nothing: NVDA scores 82.0\% on one tranche against 80.3\% on
three, and its parameter neighbourhood median \emph{improves} from 56.6\% to 58.6\%. AVGO gives
up 2.8pp of headline (90.7\% against 93.5\%) and gains on the neighbourhood, 64.8\% against
56.7\%. One tranche is both simpler and no worse.

\subsection{The rule}

Each Monday, if flat:
\begin{align*}
\text{buy} &= \min\bigl(\text{Monday open},\ \text{ATH}\times(1-\text{cap})\bigr)\\
\text{target} &= \text{buy} + \text{previous week close}\times\text{prem}
\end{align*}
where ATH is the running maximum of weekly highs through the previous week. Fill on the first
session of the week whose low reaches the buy price. Sell on any session at or after the fill
whose high reaches the target. If the target is not reached, \textbf{carry the position} and keep
the same target --- it is not recomputed while holding. Re-enter the Monday after a sale.
There is no stop loss.

\begin{table}[H]\centering\small
\begin{tabular}{lrrrrr}
\toprule
\textbf{Name} & \textbf{cap} & \textbf{prem} & \textbf{Anchor} & \textbf{Tranches} &
\textbf{Trades / yr}\\
\midrule
\rowcolor{bckeep} NVDA & 0.0935 & 0.0293 & Monday & 1 & 20\\
AVGO & 0.0800 & 0.1000 & Monday & 1 & 6.5\\
\bottomrule
\end{tabular}
\caption{Commission 0.005/share and 3.14\% interest on idle cash, as for the daily names.}
\end{table}

\subsection{On the dormant mean-reversion term}

The original weekly workbook computes an entry level
\[
L=\min\Bigl(\operatorname{mean}\bigl(m\tfrac{H_{-1}+L_{-1}}{2},\ C_{-1}w\bigr)
+\log_{10}(H_{-1}-L_{-1})\,g,\ \text{Monday open}\Bigr)
\]
with $w=1.0082$, $m=1.1627$, $g=2.0374$. On the deployed values this term \textbf{never binds}:
it sits a median 9.4\% above the Monday open and was the binding constraint in 0 of 120 NVDA
weeks and 0 of 120 AVGO weeks. The entry is therefore always the Monday open or the ATH cap, as
written in \S3.1. Retain $w$, $m$ and $g$ if replicating the sheet, but they are inert ---
$g$ can be multiplied or divided by five with no change to a single trade. Do not lower $m$:
below roughly 1.0 the term switches on and performance falls (NVDA 80.3\% to 68.5\%, AVGO 53.4\%
to 37.3\%).

\section{Allocation sheet}

\begin{table}[H]\centering\small
\begin{tabular}{lll}
\toprule
\textbf{Setting} & \textbf{Cell} & \textbf{Value}\\
\midrule
\rowcolor{bckeep} Number of names & --- & \textbf{7}\\
Weight per name & --- & \textbf{0.142857} $(=1/7)$\\
\rowcolor{bcact} Floor: min weight per stock & B8 & \textbf{0.142857}\\
Cap: max weight per stock & E11:E17 & \textbf{0.142857}\\
\rowcolor{bckeep} Bayes fraction --- flat & B6 & \textbf{0.75}\\
Split mode & --- & \textbf{1} (flat, not per-name)\\
\rowcolor{bckeep} Interest on idle cash & R2 & \textbf{0.0314}\\
\bottomrule
\end{tabular}
\caption{Equal weighting is not a default --- it was tested against inverse-volatility,
Sharpe-proportional and return-proportional schemes and beat all of them by roughly 10pp.
Setting floor and cap both to $1/7$ pins the book to equal weight.}
\end{table}

The Bayes fraction and the weight apply only to the five daily names. The two weekly names take
their $1/7$ of capital into a single Monday tranche with no sleeve split.

\section{Performance --- what to budget}

\subsection{Why the planning figures sit below the verified ones}

Three separate discounts apply, each measured rather than assumed.

\textbf{The tilt is in-sample.} The Bayes share was walk-forwarded across 30 folds. It won 20 of
them and selected a Bayes-heavy share in all 30, so the direction is supported --- but the
realised out-of-sample gain was approximately \textbf{+2pp}, against the 8--10pp the in-sample
curve suggests. Daily planning figures are therefore the 50/50 result plus 2pp, not the 75\%
result.

\textbf{The weekly parameters sit on a narrow ridge.} NVDA's 82\% is the peak of a region whose
median is 58.6\% and whose lower quartile is 45.4\%; only 22 of 169 neighbouring parameter cells
lie within 10pp of the peak. AVGO's neighbourhood median is 64.8\% with a lower quartile of
57.5\%. The neighbourhood median is what to budget, because a ridge fitted to this sample is
unlikely to stay exactly where this sample put it.

\textbf{The weekly model ends fully invested.} Every configuration finishes with the tranche
holding, so 100\% of terminal value is an open mark-to-market position rather than realised cash.
Marking it 20\% lower takes NVDA from 80.3\% to 63.6\% and AVGO from 93.5\% to 75.6\%. This is
not a bias to subtract, but it means each headline embeds one day's closing price.

\subsection{The half-sample test}

Weeks 1--59 sit inside every walk-forward training window; weeks 60--120 were scored outside at
least one. Splitting the two separates a configuration that works from one that fitted an early
run. This test is what removed four names, and the seven retained were held to the same standard.

\begin{table}[H]\centering\small
\begin{tabular}{lrrr}
\toprule
\textbf{Name} & \textbf{First half} & \textbf{Tested half} & \textbf{Shift}\\
\midrule
\rowcolor{bckeep} MU   & $-10.5$\% & 160.0\% & $+170$pp\\
VRT  & $-12.3$\% & 168.0\% & $+180$pp\\
\rowcolor{bckeep} TSM  & 40.4\%  & 88.5\%  & $+48$pp\\
VST  & 67.7\%  & 55.2\%  & $-13$pp\\
\rowcolor{bckeep} RKLB & 145.8\% & 64.7\%  & $-81$pp\\
NVDA & 44.2\%  & 81.8\%  & $+38$pp\\
\rowcolor{bckeep} AVGO & 72.2\%  & 96.4\%  & $+24$pp\\
\midrule
\rowcolor{bcwarn} \textit{Removed:} PLTR & 87.2\% & 6.6\% & $-81$pp\\
\rowcolor{bcwarn} \textit{Removed:} SPOT & 60.7\% & 3.6\% & $-57$pp\\
\rowcolor{bcwarn} \textit{Removed:} SOFI & 19.1\% & 13.0\% & $-6$pp\\
\rowcolor{bcwarn} \textit{Removed:} TSLA & $-0.2$\% & 19.2\% & $+19$pp\\
\bottomrule
\end{tabular}
\caption{All seven retained names clear 50\% in the tested half. MU and VRT are negative in the
untested half and strongest in the tested one, which is the opposite of a fitted result.}
\end{table}

\section{What was tested and rejected}

The four removed names were each put through the full weekly pipeline before being dropped ---
own-parameter grid, walk-forward, neighbourhood and half-sample split.

\begin{table}[H]\centering\small
\begin{tabular}{lrrrl}
\toprule
\textbf{Name} & \textbf{Optimised} & \textbf{Neighbourhood} & \textbf{Walk-forward} &
\textbf{Reason for removal}\\
\midrule
\rowcolor{bcwarn} PLTR & 125.9\% & 95.9\% & 1/3 & 365\% first half, 3.1\% tested half\\
SOFI & 79.7\%  & 54.8\% & 1/3 & tested half only 23.3\%\\
\rowcolor{bcwarn} TSLA & 49.4\%  & 21.0\% & 2/3 & optimum on the grid boundary; folds disagreed\\
SPOT & 35.2\%  & 24.2\% & 0/3 & tested half $-29.1$\%\\
\bottomrule
\end{tabular}
\end{table}

Also tested and rejected on the retained names: five-parameter re-optimisation of the weekly model
(0 of 3 folds), the same with $g$ frozen (0 of 3), two-parameter re-optimisation (1 of 3), and
anchor mixing across weekdays (0 of 3, and it lowers the neighbourhood median from 56.6\% to
48.2\%). Across four experiments, re-fitting the weekly parameters has won 1 fold in 9.

The Monday anchor itself \emph{is} supported. Weekend repricing is 1.75\% against 1.41\% for a
normal overnight gap, and the effect tracks the number of non-trading days rather than the
weekday label --- a Tuesday following a three-day break reprices 1.88\% against 1.35\% for an
ordinary Tuesday. Monday is the best anchor for NVDA in every sub-period and for AVGO under both
its own and NVDA's parameters.

\section{Risks}

\textbf{Concentration.} Four of the seven --- NVDA, AVGO, TSM and VRT --- are substantially the
same semiconductor and AI-infrastructure trade. The tested half was a strong period for that
complex, and the second-half strength of those four is partly one macro move counted four times.
This is the book's largest unhedged exposure and it increased when the four removed names left.

\textbf{Name-level selection.} The seven were identified as the survivors using the same
2024--26 sample on which they are scored. The half-sample test in \S5.2 was run precisely to
check this and they passed it, but a 50\% floor met by removing everything below the line is a
weaker claim than a strategy that earns 50\% per name.

\textbf{RKLB dominance.} RKLB contributes roughly 8pp of the book's 65\% planning figure. Without
it the book plans at 56\%. Its tested half is 64.7\% against a first half of 145.8\%, the largest
decline of any retained name.

\textbf{No stop on the weekly names.} A position is carried until the target is reached, however
long that takes. Median holding is 4 days for NVDA and 20 for AVGO, but the 95th percentiles are
224 and 77 days respectively.

\textbf{AVGO trade count.} A single Monday tranche produces roughly 6.5 completed trades a year
on AVGO. That is a thin base, and the name's optimum is a single grid cell that its neighbours do
not match --- 68\% at prem 0.08 and 64\% at 0.12 against 94\% at 0.10. AVGO is the least secure of
the seven.

\vfill
{\color{bcrule}\rule{\textwidth}{0.6pt}}\\[0.2em]
{\footnotesize\color{bcgrey} All returns are execution-verified: same-day round trips are
confirmed against 1-minute bars for NVDA and 5-minute bars for the remaining names, with sessions
outside minute coverage resolved only where the fill is provable at the open. Commission 0.005 per
share; idle cash earns 3.14\%.}

\end{document}
